\section{Introduction}
\label{sec:introduction}

\subsection{Large-Sample Hydrology and the CAMELS Framework}
Large-sample hydrology draws on comparative analysis across many catchments to identify controls on hydrological behavior~\citep{Addor2017}.
The CAMELS (Catchment Attributes and MEteorology for Large-sample Studies) framework has spawned regional derivatives worldwide (Table~\ref{tab:camels_comparison}), including CAMELS-GB~\citep{Coxon2020}, CAMELS-BR~\citep{Chagas2020}, CAMELS-CL~\citep{AlvarezGarreton2018}, CAMELS-AUS~\citep{Fowler2021}, LamaH-CE~\citep{Klingler2021}, CAMELS-DE~\citep{Hoege2022}, and the global Caravan compilation~\citep{Kratzert2023}.
These datasets share a standardized structure---daily discharge, meteorological forcing, and physiographic attributes---supporting hydrological model development~\citep{Kratzert2019} and prediction in ungauged basins~\citep{Wagener2004}.

\begin{table*}[htbp]
\centering
\caption{Comparison of CAMELS-RU with existing CAMELS-family datasets. Attribute counts reflect dataset-specific selections; signature counts refer to catchment-level hydrological metrics.}
\label{tab:camels_comparison}
\begin{tabular}{llrllr}
\toprule
\textbf{Dataset} & \textbf{Region} & \textbf{Catchments} & \textbf{Period} & \textbf{Attributes / Signatures} & \textbf{Area (km$^2$)} \\
\midrule
CAMELS     & USA           &    671 & 1980--2014 & 59 / 13 & 4--25{,}000 \\
CAMELS-CL  & Chile         &    516 & 1913--2016 & 70 / 13 & 5--35{,}000 \\
CAMELS-AUS & Australia     &    222 & 1951--2014 & 134 / 13 & 52--520{,}000 \\
CAMELS-GB  & Great Britain &    671 & 1970--2015 & 290 / 13 & 1--10{,}000 \\
CAMELS-BR  & Brazil        &    897 & 1980--2018 & 65 / 13 & 6--500{,}000 \\
LamaH-CE   & Central Europe &   859 & 1981--2017 & 61 / -- & 1--132{,}000 \\
CAMELS-DE  & Germany       & 1{,}582 & 1951--2020 & 69 / 13 & 1--54{,}000 \\
Caravan    & Global        & 6{,}830 & varies     & 38 / 13 & varies \\
\textbf{CAMELS-RU} & \textbf{Russia} & \textbf{\ntotal} & \textbf{\years} & \textbf{\nattributes (288) / \nsignatures} & \textbf{\minarea--\maxarea} \\
\bottomrule
\end{tabular}
\end{table*}


\subsection{The Russian Federation: An Unrepresented Region}
Although Russia spans climate zones from Arctic permafrost to semi-arid steppe, no large-sample dataset following CAMELS standards exists for the Russian Federation.
Russian hydrological data are scattered across institutional archives with non-standardized formats.
The recent closure of the AIS GMVO public data platform (2025) further underscores the need for open, preserved, and standardized hydrological records for this region.

\subsection{Objectives and Novelty}
This study presents \dataset, a quality-controlled hydrological dataset for \ntotal catchments across Russia (\years).
The dataset integrates daily discharge from \ndischarge gauging stations (AIS GMVO;~\citealt{AIS_GMVO}), meteorological forcing from ERA5-Land~\citep{Munoz2021} and MSWEP~\citep{Beck2019MSWEP}, 288 physiographic attributes from HydroATLAS~\citep{Linke2019}, and \nsignatures hydrological signatures.
Each of the \ntotal watershed boundaries was manually verified, yielding mean areal errors of \meanerror against official Roshydromet data.

Our objectives are to:
\begin{enumerate}
    \item Document data acquisition, processing, and quality control
    \item Compute hydrological signatures characterizing discharge across Russian catchments
    \item Validate meteorological forcing through inter-dataset comparison and water balance consistency
    \item Provide open-access data records following FAIR principles~\citep{Wilkinson2016}
\end{enumerate}

Sect.~\ref{sec:study_area} describes the study area;
Sect.~\ref{sec:data_sources_methods} details data sources and methods;
Sect.~\ref{sec:technical_validation} validates the dataset;
Sect.~\ref{sec:data_records} describes data records;
Sect.~\ref{sec:conclusions} summarizes contributions.
